\chapter{واکاوی ماهیت و کاربرد فرمان‌ها}

تا این مرحله، با مجموعه‌ای از دستورها آشنا شدیم که هر یک دارای گزینه‌ها (\en{Options}) و آرگومان‌های (\en{Arguments}) منحصربه‌فردی بودند. در این فصل، هدف ما ابهام‌زدایی از ماهیت این فرمان‌ها و فراتر از آن، یادگیری چگونگی تعریف دستورهای سفارشی است. دستوراتی که در این فصل مورد بررسی قرار می‌گیرند، عبارتند از:

\begin{itemize}
  \item \cmd{type}: نحوه تفسیر و ماهیت یک دستور توسط پوسته را مشخص می‌کند.
  \item \cmd{which}: مسیر مطلق فایل اجراییِ مرتبط با یک دستور را نمایش می‌دهد.
  \item \cmd{help}: راهنمای استفاده از دستورهای داخلیِ پوسته (\en{Shell Builtins}) را ارائه می‌دهد.
  \item \cmd{man}: صفحات راهنمای جامع (\en{Manual Pages}) یک دستور را نمایش می‌دهد.
  \item \cmd{apropos}: فهرستی از دستورهای مرتبط با یک کلیدواژه خاص را جستجو و ارائه می‌کند.
  \item \cmd{info}: مستندات و راهنمای یک دستور را در قالب ساختاریافته‌ی \en{Info} نمایش می‌دهد.
  \item \cmd{whatis}: شرحی بسیار کوتاه و تک‌خطی از عملکرد یک دستور ارائه می‌دهد.
  \item \cmd{alias}: یک نام مستعار یا جایگزین برای یک فرمان پیچیده ایجاد می‌کند.
\end{itemize}

\section{دسته‌بندی و ماهیت انواع فرمان‌ها در لینوکس}

یک فرمان در لینوکس می‌تواند در یکی از چهار دسته‌بندی زیر قرار گیرد:

\begin{enumerate}
  \item \textbf{برنامه‌های اجرایی (\en{Executable Programs}):} این دسته، رایج‌ترین نوع دستورها را شامل می‌شود. برنامه‌های اجرایی، فایل‌های مستقلی هستند که روی دیسک ذخیره شده‌اند. این فایل‌ها ممکن است به‌صورت باینری‌های کامپایل‌شده از زبان‌هایی مانند \en{C} یا \en{C++} باشند (نظیر دستورات \cmd{ls}، \cmd{mv} یا \cmd{cp}) و یا اسکریپت‌های تفسیری نوشته‌شده به زبان‌هایی مانند پوسته (\en{Bash})، پایتون یا پرل باشند. پوسته برای اجرای این دستورها، مسیر فایل مربوطه را در دایرکتوری‌های تعریف‌شده در متغیر محیطی \cmd{\$PATH} جستجو می‌کند.

  \item \textbf{فرمان‌های داخلی پوسته (\en{Shell Builtins}):} این دستورها مستقیماً در کد منبع خودِ پوسته (مانند \en{Bash}) گنجانده شده‌اند و به همین دلیل با سرعت و کارایی بسیار بالایی اجرا می‌شوند. فرمان‌های داخلی نیازی به فراخوانی فایل خارجی ندارند، زیرا بخشی از خودِ فرآیند پوسته محسوب می‌شوند. دستوراتی نظیر \cmd{cd} (تغییر دایرکتوری)، \cmd{pwd} (نمایش مسیر فعلی) و \cmd{echo} (چاپ متن) از این دسته‌اند. از آنجا که این فرمان‌ها به پوسته وابسته‌اند، رفتار آن‌ها ممکن است در پوسته‌های مختلف (مانند \en{Bash}، \en{Zsh} یا \en{Fish}) تفاوت‌های اندکی داشته باشد.

  \item \textbf{توابع پوسته (\en{Shell Functions}):} توابع پوسته، قطعه‌کدهای کوچکی هستند که در قالب اسکریپت‌های ساده تعریف شده و در حافظه بارگذاری می‌شوند. این توابع به کاربر اجازه می‌دهند مجموعه‌ای از دستورات را تحت یک نام واحد بسته‌بندی و فراخوانی کند. توابع معمولاً در فایل‌های پیکربندی نظیر \file{\textasciitilde/.bashrc} تعریف می‌شوند تا در ابتدای هر نشست (\en{Session})، آماده استفاده باشند. بهره‌گیری از توابع به سازمان‌دهی منطقی دستورهای پیچیده و افزایش خوانایی اسکریپت‌ها کمک شایانی می‌کند.

  \item \textbf{نام‌های مستعار (\en{Aliases}):} نام‌های مستعار یا \en{Aliases}، میان‌برهای سفارشی‌سازی‌شده‌ای هستند که کاربر برای یک یا ترکیبی از چند دستور ایجاد می‌کند. هدف اصلی نام‌های مستعار، ساده‌سازی فرمان‌های طولانی یا افزودن گزینه‌های ایمنی (مانند الحاق گزینه \cmd{-i} به دستور \cmd{rm}) است. برای نمونه، تعریف نام مستعاری به شرح زیر:

\begin{latin}
\begin{lstlisting}
ls='ls --color=auto'
\end{lstlisting}
\end{latin}

موجب می‌شود خروجی دستور \cmd{ls} همواره رنگی نمایش داده شود. این نام‌های مستعار اغلب در فایل‌های پیکربندی پوسته (مانند \file{\textasciitilde/.bash\_profile}) تعریف می‌شوند تا در تمامی نشست‌ها به‌صورت پایدار در دسترس باشند.
\end{enumerate}

\section{شناسایی ماهیت فرمان‌ها}

سیستم‌عامل لینوکس برای تشخیص دقیق اینکه یک فرمان در کدام‌یک از دسته‌بندی‌های چهارگانه فوق قرار می‌گیرد، ابزارهای تخصصی متنوعی را در اختیار کاربر قرار داده است.

\subsection{دستور type: تعیین ماهیت و نحوه تفسیر فرمان‌ها}

دستور \cmd{type} خود یک فرمان داخلی پوسته (\en{Shell Builtin}) محسوب می‌شود که وظیفه آن نمایش نحوه تفسیر یک دستور توسط پوسته است. این ابزار کلیدی به شما اجازه می‌دهد تا به‌سرعت تشخیص دهید که آیا یک فرمان یک برنامه اجرایی مستقل است، یا در قالب یک دستور داخلی، تابع و یا نام مستعار (\en{Alias}) عمل می‌کند.

نحو (\en{Syntax}) کلی این دستور به شرح زیر است:

\begin{latin}
\begin{lstlisting}
type command
\end{lstlisting}
\end{latin}

در این ساختار، \file{command} نام دستوری است که قصد واکاوی ماهیت آن را دارید. در ادامه، با تحلیل چند مثال کاربردی، عملکرد دقیق این ابزار را بررسی می‌کنیم:

\begin{latin}
\begin{lstlisting}
[me@linuxbox ~]$ type type
type is a shell builtin
[me@linuxbox ~]$ type ls
ls is aliased to `ls --color=auto'
[me@linuxbox ~]$ type cp
cp is /usr/bin/cp
\end{lstlisting}
\end{latin}

\textbf{تحلیل فنی خروجی‌های دستور \cmd{type}:}

\begin{itemize}
  \item \textbf{بررسی \cmd{type type}:} خروجی \en{type is a shell builtin} بیانگر این است که خودِ دستور \cmd{type} نیز یک فرمان داخلی پوسته است؛ به این معنا که به‌صورت مستقیم در کد منبع پوسته گنجانده شده و فاقد فایل اجرایی مجزا روی دیسک است.

  \item \textbf{بررسی \cmd{type ls}:} خروجی \en{'ls is aliased to 'ls --color=auto} یک واقعیت فنی مهم را آشکار می‌سازد؛ در بسیاری از توزیع‌های مدرن لینوکس، آنچه کاربر به‌عنوان دستور \cmd{ls} فراخوانی می‌کند، در حقیقت یک نام مستعار (\en{Alias}) است که به‌طور خودکار گزینه \cmd{--color=auto} را به دستور اصلی الحاق می‌کند. علت رنگی بودن خروجی لیست فایل‌ها در ترمینال نیز همین موضوع است.

  \item \textbf{بررسی \cmd{type cp}:} خروجی \en{cp is /usr/bin/cp} تایید می‌کند که \cmd{cp} یک برنامه اجرایی مستقل (\en{Executable}) است. در این حالت، دستور \cmd{type} مسیر مطلق فایل در سلسله‌مراتب سیستم فایل را نیز گزارش می‌دهد.
\end{itemize}

\subsection{دستور which: شناسایی مسیر برنامه‌های اجرایی}

در سناریوهای خاص، به‌ویژه در محیط‌های سروری یا سیستم‌های توسعه، ممکن است چندین نسخه از یک برنامه اجرایی به‌صورت هم‌زمان نصب شده باشند. دستور \cmd{which} ابزاری تخصصی برای شناسایی موقعیت دقیق فایل اجرایی است که پوسته در زمان فراخوانی یک دستور، آن را اجرا می‌کند. این دستور با پیمایش دایرکتوری‌های تعریف‌شده در متغیر محیطی \cmd{\$PATH} (که اولویت‌بندی جستجوی دستورها را تعیین می‌کند)، نخستین مسیر منطبق را یافته و خروجی می‌دهد.

\textbf{مثال عملی:}

\begin{latin}
\begin{lstlisting}
[me@linuxbox ~]$ which ls
/usr/bin/ls
\end{lstlisting}
\end{latin}

در این نمونه، خروجی مشخص می‌کند که نسخه اصلی و قابل‌اجرای فرمان \cmd{ls} در مسیر \file{/usr/bin/} مستقر است.

فرمان \cmd{which} منحصراً برای ردیابی برنامه‌های اجرایی (\en{Executable Programs}) طراحی شده است. این ابزار فاقد توانایی لازم برای شناسایی و تحلیل دستورهای داخلی پوسته (\en{Shell Builtins}) یا نام‌های مستعار (\en{Aliases}) است. چنانچه این دستور را برای واکاوی یک فرمان داخلی مانند \cmd{cd} به‌کار بگیرید، با پاسخی مبنی بر عدم یافتن فایل مواجه خواهید شد:

\begin{latin}
\begin{lstlisting}
[me@linuxbox ~]$ which cd
/usr/bin/which: no cd in (/usr/local/bin:/usr/bin:/bin:/usr/local
/games:/usr/games)
\end{lstlisting}
\end{latin}

این خروجی به‌روشنی تایید می‌کند که \cmd{cd} به‌عنوان یک فایل مستقل در مسیرهای تعریف‌شده در متغیر محیطی \cmd{\$PATH} وجود ندارد، چرا که این فرمان به‌صورت درونی در ساختار خودِ پوسته تعبیه شده است. از این رو، در دنیای لینوکس برای دستیابی به یک تحلیل جامع و دقیق از ماهیت هر نوع دستور، فرمان \cmd{type} ابزاری به‌مراتب کارآمدتر و قابل‌اطمینان‌تر محسوب می‌شود.

\section{روش‌های دسترسی به مستندات و راهنمای دستورها}

پس از شناسایی ماهیت فنی یک فرمان، گام بعدی دستیابی به مستندات رسمی آن است. این مستندات نقشه راه دقیقی برای درک نحوه به‌کارگیری (\en{Syntax})، گزینه‌ها (\en{Options}) و رفتار عملیاتی هر دستور در اختیار کاربر قرار می‌دهند.

\subsection{دستور help: مستندات اختصاصی فرمان‌های داخلی پوسته}

فرمان \cmd{help} یک ابزار درونی در پوسته \en{Bash} است که منحصراً برای ارائه مستندات دستورهای داخلی (\en{Shell Builtins}) طراحی شده است. برای استفاده از این ابزار، کافی است واژه \cmd{help} را پیش از نام دستور مورد نظر درج کنید.

\textbf{مثال عملی:}

\begin{latin}
\begin{lstlisting}
[me@linuxbox ~]$ help cd
cd: cd [-L|[-P [-e]] [-@]] [dir]
    Change the shell working directory.

    Change the current directory to DIR.  The default DIR is the
    value of the HOME shell variable.

    The variable CDPATH defines the search path for the directory
    containing DIR.  Alternative directory names in CDPATH are
    separated by a colon (:). A null directory name is the same as
    the current directory.  If DIR begins with a slash (/), then
    CDPATH is not used.

    If the directory is not found, and the shell option `cdable_vars'
    is set, the word is assumed to be  a variable name.  If that
    variable has a value, its value is used for DIR.

    Options:
      -L    force symbolic links to be followed: resolve symbolic
            links in DIR after processing instances of `..'.
      -P    use the physical directory structure without following
            symbolic links: resolve symbolic links in DIR before
            processing instances of `..'.
      -e    if the -P option is supplied, and the current working
            directory cannot be determined successfully, exit with
            a non-zero status
      -@    on systems that support it, present a file with extended
            attributes as a directory containing the file attributes

    The default is to follow symbolic links, as if `-L' were
    specified. `..' is processed by removing the immediately previous
    pathname component back to a slash or the beginning of DIR.

    Exit Status:
    Returns 0 if the directory is changed, and if $PWD is set
    successfully when -P is used; non-zero otherwise.
\end{lstlisting}
\end{latin}

خروجی این دستور، راهنمای جامع \cmd{cd} را به‌صورت بالا نمایش می‌دهد. این خروجی، برخلاف صفحات راهنمای حجیم، بر روی تعامل مستقیم دستور با محیط پوسته تمرکز دارد و سریع‌ترین راه برای یادآوری گزینه‌های یک فرمان داخلی است.

برای درک صحیح و دقیق متون راهنما، تسلط بر نشانه‌گذاری‌های استاندارد در دنیای لینوکس ضروری است:

\begin{itemize}
  \item \textbf{براکت‌ها \cmd{[]}:} هر عبارت یا گزینه‌ای که درون براکت قرار می‌گیرد، «اختیاری» (\en{Optional}) است و عدم استفاده از آن در اجرای دستور خللی ایجاد نمی‌کند.
  \item \textbf{خط عمودی \cmd{|}:} این نماد بیانگر «انتخاب یکی از گزینه‌ها» است؛ به این معنا که در یک لحظه تنها می‌توان یکی از گزینه‌های جدا شده توسط این علامت را برگزید.
\end{itemize}

با واکاوی نحو (\en{Syntax}) دستور \cmd{cd} در مثال زیر:

\begin{latin}
\begin{lstlisting}
cd [-L | [-P [-e]]] [dir]
|   |      |    |       |
|   |      |    |       +-- dir: optional
|   |      |    +---------- -e: optional (only with -P)
|   |      +--------------- -P: optional
|   +---------------------- -L: optional (alternative to -P)
+-------------------------- cd: the command itself
\end{lstlisting}
\end{latin}

\begin{itemize}
  \item استفاده از گزینه‌های \cmd{-L} یا \cmd{-P} کاملاً اختیاری است.
  \item به دلیل وجود خط عمودی، استفاده هم‌زمان از \cmd{-L} و \cmd{-P} مجاز نیست و باید یکی را انتخاب کرد.
  \item در صورت انتخاب گزینه \cmd{-P} (و تنها در این صورت)، می‌توان به‌صورت اختیاری از گزینه‌ی \cmd{-e} نیز استفاده نمود.
  \item پارامتر مربوط به دایرکتوری (\file{dir}) نیز در انتهای دستور اختیاری در نظر گرفته شده است.
\end{itemize}

خروجی حاصل از فرمان \cmd{help} معمولاً بسیار فشرده و تخصصی است؛ لذا طبیعی است که برخی مفاهیم فنی آن در ابتدای مسیر یادگیری برای شما مبهم به نظر برسد. این پیچیدگی با تداوم مطالعه و تجربه عملی مرتفع خواهد شد.

\begin{note}
\textbf{نکته:} با به‌کارگیری گزینه \cmd{-m} در کنار دستور \cmd{help} (به‌صورت \cmd{help -m command})، خروجی راهنما با قالب‌بندی متفاوتی نمایش می‌یابد که به استانداردهای صفحات راهنما (\en{Man Pages}) نزدیک‌تر بوده و خوانایی به‌مراتب بیشتری دارد.
\end{note}

\subsection{گزینه‌ی --help: نمایش راهنمای کاربری و نحوه‌ی اجرای دستور}

بسیاری از برنامه‌های اجرایی در لینوکس از گزینه \cmd{--help} پشتیبانی می‌کنند. این گزینه، خلاصه‌ای کاربردی از نحو نگارش (\en{Syntax})، گزینه‌ها (\en{Options}) و آرگومان‌های ورودی دستور را در اختیار کاربر قرار می‌دهد. استفاده از این قابلیت، سریع‌ترین روش برای دستیابی به اطلاعات عملیاتی یک برنامه، بدون نیاز به پیمایش در صفحات مفصل راهنما است.

به‌عنوان نمونه، با اجرای فرمان \cmd{mkdir --help} خروجی زیر نمایش می‌یابد:

\begin{latin}
\begin{lstlisting}
[me@linuxbox ~]$ mkdir --help
Usage: mkdir [OPTION] DIRECTORY...
Create the DIRECTORY(ies), if they do not already exist.

  -Z, --context=CONTEXT (SELinux) set security context to CONTEXT
 Mandatory arguments to long options are mandatory for short options too.
  -m, --mode=MODE    set file mode (as in chmod), not a=rwx - umask
  -p, --parents      no error if existing, make parent directories as
                     needed
  -v, --verbose      print a message for each created directory
      --help         display this help and exit
      --version      output version information and exit
 Report bugs to <bug-coreutils@gnu.org>.
\end{lstlisting}
\end{latin}

همان‌طور که ملاحظه می‌شود، این خروجی ساختار کلی فراخوانی دستور (\en{Usage}) و شرح مختصری از عملکرد هر یک از گزینه‌ها را ارائه می‌دهد. همچنین، نمایش هم‌زمان گزینه‌های کوتاه (مانند \cmd{-m}) و معادل‌های بلند آن‌ها (مانند \cmd{--mode}) فرآیند یادگیری و به‌خاطرسپاری دستورات را تسهیل می‌کند.

شایان ذکر است که برخی از برنامه‌ها ممکن است از قرارداد \cmd{--help} پیروی نکنند. در چنین مواردی، استفاده از این گزینه معمولاً منجر به صدور یک پیام خطا از سوی برنامه می‌شود؛ با این حال، بسیاری از برنامه‌ها در همان پیام خطا نیز اطلاعات زیربنایی مربوط به نحوِ صحیحِ اجرای دستور را به کاربر ارائه می‌دهند که می‌تواند بسیار راهگشا باشد.

\subsection{دستور man: نمایش صفحات راهنمای جامع (Manual Pages)}

تقریباً تمامی برنامه‌های اجرایی که برای تعامل در محیط خط فرمان لینوکس طراحی شده‌اند، دارای مستندات رسمی و مرجعی به نام صفحه راهنما یا \en{man page} هستند. برای دسترسی به این گنجینه اطلاعاتی، از ابزار اختصاصی \cmd{man} استفاده می‌شود. نحوه‌ی فراخوانی آن بسیار ساده و به صورت زیر است:

\begin{latin}
\begin{lstlisting}
man program
\end{lstlisting}
\end{latin}

صفحات \cmd{man} به صورت استاندارد از بخش‌های زیر تشکیل شده‌اند:

\begin{itemize}
  \item \textbf{عنوان (\en{NAME}):} نام دقیق دستور و شرحی بسیار کوتاه از عملکرد آن.
  \item \textbf{خلاصه نحو (\en{SYNOPSIS}):} الگوی فنی فراخوانی دستور، شامل تمامی گزینه‌ها و آرگومان‌های مجاز.
  \item \textbf{توضیحات (\en{DESCRIPTION}):} تشریح جامع و مفصل اهداف، رفتارها و منطق عملیاتی دستور.
  \item \textbf{گزینه‌ها (\en{OPTIONS}):} فهرست کامل گزینه‌ها به همراه جزئیات فنی مربوط به هر یک.
\end{itemize}

باید توجه داشت که صفحات \cmd{man} ماهیتی کاملاً مرجع‌گونه (\en{Reference}) دارند و هدف آن‌ها ارائه دقیق‌ترین جزئیات فنی است؛ لذا به ندرت شامل مثال‌های آموزشی یا سناریوهای کاربردی هستند. به همین دلیل، این مستندات ممکن است برای کاربران تازه‌وارد کمی دشوار به نظر برسند. برای نمونه، جهت مطالعه مستندات کامل فرمان \cmd{ls} می‌توان دستور زیر را اجرا کرد:

\begin{latin}
\begin{lstlisting}
[me@linuxbox ~]$ man ls
\end{lstlisting}
\end{latin}

در اکثر توزیع‌های لینوکس، ابزار \cmd{man} برای نمایش محتوا از برنامه‌ی جانبی \cmd{less} استفاده می‌کند. به همین سبب، تمامی کلیدهای کنترلی و قابلیت‌های پیمایشی \cmd{less} در هنگام مطالعه‌ی صفحات راهنما نیز در دسترس هستند. برخی از پرکاربردترینِ این دستورات عبارتند از:

\begin{itemize}
  \item \cmd{G} (حرف بزرگ): انتقال آنی به انتهای صفحه.
  \item \cmd{g} (حرف کوچک): بازگشت سریع به ابتدای صفحه.
  \item \cmd{Space} (کلید فاصله): پیمایش به سمت پایین به اندازه‌ی یک صفحه.
  \item \cmd{Enter}: پیمایش خط‌به‌خط به سمت پایین.
  \item \cmd{q} (خروج): بستن صفحه راهنما و بازگشت به محیط پوسته.
\end{itemize}

علاوه بر این موارد، شما می‌توانید با فشردن کلید \cmd{/} و سپس وارد کردن یک عبارت، در میان متنِ راهنما به جستجوی کلیدواژه‌های خاص بپردازید تا به بخش مورد نظر خود سریع‌تر دسترسی پیدا کنید.

مستندات ارائه‌شده توسط سیستم \cmd{man} به بخش‌های مجزایی تقسیم می‌شوند که نه‌تنها فرمان‌های کاربری، بلکه دستورالعمل‌های مدیریتی، رابط‌های برنامه‌نویسی (\en{API})، ساختار فایل‌های پیکربندی و موارد دیگر را پوشش می‌دهند. این تقسیم‌بندی به کاربر کمک می‌کند تا در صورت وجود نام‌های مشابه، مستندات دقیق مورد نیاز خود را بیابد.

\begin{table}[htbp]
\centering
\caption{ساختار و طبقه‌بندی بخش‌های صفحات \cmd{man}}
\label{tab:man-sections}
\begin{tabularx}{\textwidth}{c X}
\toprule
\textbf{شماره بخش} & \textbf{محتوا} \\
\midrule
\cmd{1} & فرمان‌های کاربری (\en{User Commands}): شامل دستوراتی که کاربران عادی به‌طور روزمره با آن‌ها تعامل دارند، مانند \cmd{ls}، \cmd{cp}، \cmd{mv} و \cmd{rm}. \\
\addlinespace
\cmd{2} & فراخوان‌های سیستمی هسته (\en{Kernel System Calls}): توابعی که به‌عنوان رابط میان برنامه‌ها و هسته سیستم‌عامل عمل می‌کنند. \\
\addlinespace
\cmd{3} & توابع کتابخانه استاندارد \en{C}: مستندات مربوط به کتابخانه استاندارد \en{C} که برنامه‌نویسان در توسعه نرم‌افزار از آن‌ها بهره می‌برند. \\
\addlinespace
\cmd{4} & فایل‌های ویژه (\en{Special Files}): فایل‌های مرتبط با گره‌های دستگاه (\en{Device Nodes}) و درایورهای سیستم. \\
\addlinespace
\cmd{5} & قالب‌های فایل (\en{File Formats}): شرح ساختار و نحوِ نگارش فایل‌های پیکربندی سیستم، نظیر \file{/etc/passwd} یا \file{/etc/fstab}. \\
\addlinespace
\cmd{6} & بازی‌ها و سرگرمی‌ها (\en{Games \& Amusements}): مستندات مربوط به برنامه‌های تفریحی و بازی‌های تحت خط فرمان. \\
\addlinespace
\cmd{7} & مباحث متفرقه (\en{Miscellaneous}): شامل استانداردهای سیستم، پروتکل‌های شبکه و مفاهیم کلی که در دسته‌های دیگر نمی‌گنجند. \\
\addlinespace
\cmd{8} & مدیریت سیستم (\en{System Administration Commands}): فرمان‌های سطح بالا که مخصوص مدیران سیستم (\en{Root}) برای پیکربندی و نگهداری است، مانند \cmd{mount} و \cmd{reboot}. \\
\bottomrule
\end{tabularx}
\end{table}

در مواردی که یک نام واحد (نظیر \file{passwd}) به‌طور هم‌زمان هم به یک فرمان اجرایی و هم به یک مفهوم دیگر (مانند فایل پیکربندی) اطلاق شود، جهت دستیابی به مستندات هدفمند، تعیین شماره بخش (\en{Section Number}) الزامی است. در صورت عدم درج شماره بخش، ابزار \cmd{man} به‌صورت پیش‌فرض نخستین مورد انطباق‌یافته در سلسله‌مراتب خود را که معمولاً در بخش ۱ (فرمان‌های کاربری) قرار دارد، نمایش می‌دهد.

برای فراخوانی مستندات از یک بخش خاص، از نحوِ (\en{Syntax}) زیر استفاده می‌شود:

\begin{latin}
\begin{lstlisting}
man section search_term
\end{lstlisting}
\end{latin}

در این الگو، \file{section} نشان‌دهنده شماره بخش مورد نظر و \file{search\_term} بیانگر موضوع یا فرمانی است که قصد واکاوی آن را دارید.

\textbf{نمونه:} جهت مطالعه مستندات فنی مربوط به ساختار و قالب فایل \file{/etc/passwd} که در بخش ۵ (\en{File Formats}) طبقه‌بندی شده است، فرمان زیر را اجرا می‌کنیم:

\begin{latin}
\begin{lstlisting}
[me@linuxbox ~]$ man 5 passwd
\end{lstlisting}
\end{latin}

با اجرای این دستور، \cmd{man} به‌جای نمایش راهنمای کاربردیِ دستورِ تغییر کلمه عبور (بخش ۱)، مستندات تخصصی مربوط به ساختارِ داده‌ای فایلِ پسورد را ارائه می‌دهد. این رویکرد، ابزاری دقیق برای تفکیک مفاهیم هم‌نام و دسترسی به جزئیات فنی در سطوح مختلف سیستم‌عامل فراهم می‌آورد.

\subsection{دستور apropos: جستجوی موضوعی در صفحات راهنما}

در بسیاری از سناریوها، ممکن است نام دقیق یک فرمان را به خاطر نیاورید، اما با عملکرد یا کلیدواژه‌های مرتبط با آن آشنا باشید. ابزار \cmd{apropos} به شما امکان می‌دهد تا در میان عناوین و شرح کوتاهِ تمامی صفحات راهنما (\en{man pages}) به جستجوی یک عبارت خاص بپردازید. این قابلیت برای یافتن ابزارهای ناشناخته یا جایگزین بسیار حیاتی است.

در نمونه زیر، با جستجوی واژه \en{``partition''}، فهرستی از تمامی ابزارهای مرتبط با مدیریت پارتیشن‌های دیسک نمایش داده می‌شود:

\begin{latin}
\begin{lstlisting}
[me@linuxbox ~]$ apropos partition
addpart (8)          - simple wrapper around the "add partition"...
all-swaps (7)        - event signalling that all swap partitions...
cfdisk (8)           - display or manipulate disk partition table
cgdisk (8)           - Curses-based GUID partition table (GPT)...
delpart (8)          - simple wrapper around the "del partition"...
fdisk (8)            - manipulate disk partition table
fixparts (8)         - MBR partition table repair utility
gdisk (8)            - Interactive GUID partition table (GPT)...
mpartition (1)       - partition an MSDOS hard disk
partprobe (8)        - inform the OS of partition table changes
partx (8)            - tell the Linux kernel about the presence...
resizepart (8)       - simple wrapper around the "resize partition"...
sfdisk (8)           - partition table manipulator for Linux
sgdisk (8)           - Command-line GUID partition table (GPT)...
\end{lstlisting}
\end{latin}

هر سطر از خروجی شامل دو بخش کلیدی است: بخش نخست، نامِ صفحه راهنما و بخش دوم (که در پرانتز درج شده)، شماره بخش (\en{Section}) مربوطه را نمایش می‌دهد. این اطلاعات به شما اجازه می‌دهد تا با استفاده از فرمان \cmd{man} و شماره بخش ارائه شده، مستندات دقیق آن ابزار را فراخوانی کنید.

شایان ذکر است که فرمان \cmd{man -k} از نظر عملکردی کاملاً با \cmd{apropos} منطبق است و در ادبیات سیستم‌عامل لینوکس، این دو دستور اغلب به‌صورت مترادف به‌کار می‌روند.

\subsection{دستور whatis: نمایش خلاصه توصیف عملکرد فرمان}

فرمان \cmd{whatis} ابزاری کارآمد برای نمایش نام و شرحی بسیار کوتاه از صفحات راهنما است که با کلیدواژه‌ی مدنظر تطابق دارند. این ابزار در واقع خلاصه‌ای کوتاه از کارکرد یک دستور را ارائه می‌دهد و زمانی که تنها به یک یادآوری سریع از ماهیت فنی یک فرمان نیاز دارید، جایگزینی سریع برای مطالعه‌ی کامل صفحات \cmd{man} محسوب می‌شود.

\textbf{مثال عملی:}

\begin{latin}
\begin{lstlisting}
[me@linuxbox ~]$ whatis ls
ls                   (1)  - list directory contents
\end{lstlisting}
\end{latin}

در این نتیجه، \cmd{ls} نمایانگر نام فرمان و عدد \cmd{(1)} بیانگر طبقه‌بندی آن در بخش اول (فرمان‌های کاربری) صفحات راهنما است. در ادامه‌ی سطر، یک توصیفِ کوتاه و دقیق از عملکرد اصلی دستور (در اینجا: فهرست کردن محتویات دایرکتوری) درج شده است.

\section{تحلیل تخصصی صفحه راهنمای پوسته bash}

صفحات راهنمای لینوکس و سیستم‌های شبه‌یونیکس ماهیتی صرفاً مرجع (\en{Reference}) دارند و به هیچ وجه به‌عنوان متون آموزشی تدوین نشده‌اند. در میان این مستندات، جایزه دشوارترین صفحه راهنما احتمالاً به \cmd{bash} تعلق می‌گیرد. این سند به‌طرز حیرت‌انگیزی متراکم، فنی و لبریز از جزئیات است؛ به‌گونه‌ای که نسخه چاپی آن بالغ بر ۸۰ صفحه برآورد می‌شود. ساختار این راهنما برای کاربران تازه‌کار بسیار پیچیده و فراتر از سطح درک اولیه است، اما همین دقت و جامعیت، آن را به معتبرترین منبع برای متخصصان تبدیل کرده است. بررسی این صفحه می‌تواند چالشی جدی برای سنجش میزان تسلط شما بر مفاهیم عمیق پوسته باشد.

\section{دستور info: نمایش و پیمایش مستندات گنو (GNU)}

پروژه \en{GNU} با هدف ارائه جایگزینی مدرن برای صفحات سنتی \cmd{man} (که محدود به متن ساده بودند)، مستندات خود را در قالب ساختاری به نام \cmd{info} عرضه کرد. این مستندات توسط برنامه‌ای با همین نام مدیریت می‌شوند و با بهره‌گیری از مفاهیمی چون ابرپیوندها (\en{Hyperlinks})، شباهت ساختاری زیادی به صفحات وب دارند که پیمایش میان بخش‌های مختلف را بسیار تسهیل می‌کند.

ابزار \cmd{info} فایل‌های اطلاعاتی را در یک ساختار درختی و در قالب گره‌های مجزا (\en{Nodes}) سازمان‌دهی می‌کند. هر گره به یک سرفصل اختصاصی می‌پردازد و می‌تواند حاوی پیوندهایی به گره‌های دیگر باشد. این پیوندها با علامت ستاره (\cmd{*}) در ابتدای خط متمایز می‌شوند؛ برای فعال‌سازی آن‌ها کافی است مکان‌نما را روی پیوند قرار داده و کلید \en{Enter} را بفشارید.

در محیط \cmd{info} (مانند نمونه زیر)، اطلاعات در گره‌های تخصصی سازمان‌دهی شده‌اند که در بالای صفحه، موقعیت گره فعلی نسبت به گره‌های بعدی (\en{Next}) و گره بالادستی یا والد (\en{Up}) مشخص شده است:

\begin{latin}
\begin{lstlisting}
File: coreutils.info,  Node: ls invocation,  Next: dir invocation,
Up: Directory listing
10.1 `ls': List directory contents
=================================
The `ls' program lists information about files (of any type,
including directories). Options and file arguments can be intermixed
arbitrarily, as usual.
  For non-option command-line arguments that are directories, by
default `ls' lists the contents of directories, not recursively, and
omitting files with names beginning with `.'. For other non-option
arguments, by default `ls' lists just the filename. If no non-option
argument is specified, `ls' operates on the current directory, acting
as if it had been invoked with a single argument of `.'.

  By default, the output is sorted alphabetically, according to the
--zz-Info: (coreutils.info.gz)ls invocation, 63 lines --Top----------
\end{lstlisting}
\end{latin}

این مستندات توضیح می‌دهند که برنامه \cmd{ls} به‌صورت پیش‌فرض محتویات دایرکتوری‌ها را فهرست می‌کند (بدون پیمایش بازگشتی و با صرف‌نظر از فایل‌های پنهان). در صورت عدم درج آرگومان، دستور روی دایرکتوری جاری (\cmd{.}) اعمال می‌شود و خروجی نیز به‌صورت الفبایی مرتب می‌گردد.

برای فراخوانی این محیط تعاملی، کافی است نام ابزار \cmd{info} را پیش از نام برنامه مورد نظر درج کنید:

\begin{latin}
\begin{lstlisting}
info ls
\end{lstlisting}
\end{latin}

\begin{table}[htbp]
\centering
\caption{کلیدهای کنترلی و فرمان‌های پیمایش محیط \cmd{info}}
\label{tab:info-keys}
\begin{tabularx}{\textwidth}{l X}
\toprule
\textbf{کلید میان‌بر} & \textbf{عملکرد و نتیجه عملیاتی} \\
\midrule
\cmd{?} & فراخوانی راهنمای جامع تمامی دستورات و کلیدهای پیمایشی محیط \cmd{info}. \\
\addlinespace
\en{PgUp} یا \en{Backspace} & بازگشت به صفحه قبلی در گره جاری. \\
\addlinespace
\en{PgDn} یا \cmd{Space} & انتقال به صفحه بعدی در گره جاری. \\
\addlinespace
\cmd{n} & \en{Next}: انتقال به گره بعدی در همان سطح از ساختار درختی مستندات. \\
\addlinespace
\cmd{p} & \en{Previous}: بازگشت به گره قبلی در همان سطح از ساختار درختی. \\
\addlinespace
\cmd{u} & \en{Up}: صعود به گره والد (سطح بالاتر) که معمولاً دربرگیرنده فهرست اصلی یا منوی فوقانی است. \\
\addlinespace
\en{Enter} & فعال‌سازی و دنبال کردن ابرپیوند (\en{Hyperlink}) که مکان‌نما (\en{Cursor}) بر روی آن قرار دارد. \\
\addlinespace
\cmd{q} & خروج از محیط برنامه \cmd{info} و بازگشت به محیط پوسته (\en{Shell}). \\
\bottomrule
\end{tabularx}
\end{table}

بسیاری از ابزارهای خط فرمانی که تاکنون بررسی کرده‌ایم، زیرمجموعه‌ای از بسته‌ی نرم‌افزاری \en{coreutils} در پروژه‌ی \en{GNU} محسوب می‌شوند. برای دسترسی به مستندات جامع و یکپارچه‌ی این مجموعه، می‌توانید فرمان زیر را اجرا کنید:

\begin{latin}
\begin{lstlisting}
[me@linuxbox ~]$ info coreutils
\end{lstlisting}
\end{latin}

اجرای این دستور، یک صفحه‌ی اطلاعاتی مرجع به‌همراه یک منوی درختی را فراخوانی می‌کند. این منو شامل ابرپیوندها (\en{Hyperlinks}) به مستندات اختصاصی هر یک از برنامه‌های موجود در بسته‌ی \en{coreutils} است و بستری ایده‌آل برای پیمایش سریع و ساختاریافته میان ابزارهای مختلف فراهم می‌سازد.

\section{بازیابی مستندات محلی و فایل‌های راهنما در /usr/share/doc/}

بسیاری از بسته‌های نرم‌افزاری نصب‌شده بر روی سیستم، فایل‌های مستندات اختصاصی خود را در مسیر استاندارد \file{/usr/share/doc/} نگهداری می‌کنند. این دایرکتوری مرجعی ارزشمند برای درک جزئیات پیاده‌سازی و پیکربندی هر نرم‌افزار است. این فایل‌ها معمولاً در قالب‌های زیر ارائه می‌شوند:

\begin{itemize}
  \item \textbf{متن ساده (\en{Plain Text}):} بخش عمده‌ای از این مستندات به صورت متن خام هستند که با استفاده از ابزارهای نمایش متن (\en{Pagers}) مانند \cmd{less} به‌راحتی قابل مطالعه‌اند.

  \item \textbf{قالب \en{HTML}:} برخی بسته‌ها مستندات خود را در قالب صفحات وب ارائه می‌دهند که برای مشاهده‌ی بهینه‌ی آن‌ها می‌توان از مرورگرهای وب استفاده کرد.

  \item \textbf{فایل‌های فشرده:} در بسیاری از موارد، فایل‌های مستندات با پسوند \file{.gz} ذخیره می‌شوند که نشان‌دهنده‌ی فشرده‌سازی آن‌ها توسط ابزار \cmd{gzip} جهت بهینه‌سازی فضای ذخیره‌سازی است. برای مطالعه‌ی این فایل‌ها نیازی به استخراج دستی (\en{Decompress}) آن‌ها نیست؛ ابزار \cmd{zless} به‌طور اختصاصی برای این هدف طراحی شده است. \cmd{zless} در واقع نسخه‌ای از ابزار \cmd{less} است که محتوای فایل‌های فشرده‌شده با \cmd{gzip} را به‌صورت آنی باز کرده و به کاربر نمایش می‌دهد، بدون آنکه محتوای فایل روی دیسک تغییری کند.
\end{itemize}

\section{تعریف دستورات سفارشی با نام‌های مستعار (alias)}

ایجاد دستورات سفارشی، نخستین گام در مسیر تسلط بر اسکریپت‌نویسی پوسته (\en{Shell Scripting}) محسوب می‌شود. این قابلیت از طریق فرمان \cmd{alias} در دسترس است. پیش از تعریف یک نام مستعار، آموختن یک تکنیک کلیدی ضروری است: قابلیت اجرای متوالی چندین دستور در یک خط واحد با استفاده از نویسه‌ی سمی‌کالن (\cmd{;}) به عنوان جداکننده.

\begin{latin}
\begin{lstlisting}
command1; command2; command3...
\end{lstlisting}
\end{latin}

به نمونه‌ی عملیاتی زیر توجه کنید:

\begin{latin}
\begin{lstlisting}
[me@linuxbox ~]$ cd /usr; ls; cd -
bin  games  include  lib  local  sbin  share  src
/home/me
[me@linuxbox ~]$
\end{lstlisting}
\end{latin}

در این مثال، زنجیره‌ای از فرامین به‌صورت متوالی اجرا می‌شوند: ابتدا محیط کاری به دایرکتوری \file{/usr/} تغییر می‌یابد، سپس محتویات آن فهرست (\cmd{ls}) شده و در نهایت با استفاده از آرگومان \cmd{-} در دستور \cmd{cd} (معادل دایرکتوری کاری قبلی)، کاربر به موقعیت اولیه‌ی خود بازمی‌گردد.

برای ایجاد یک فرمان سفارشی، نخستین گام انتخاب نامی است که با دستورات سیستمی موجود تداخل نداشته باشد. پیش از نهایی کردن نام، توصیه می‌شود با استفاده از ابزار \cmd{type} وضعیت آزاد بودن آن را بررسی کنید:

\begin{latin}
\begin{lstlisting}
[me@linuxbox ~]$ type test
test is a shell builtin
\end{lstlisting}
\end{latin}

در نمونه‌ی فوق، ملاحظه می‌کنید که نام \cmd{test} پیش‌تر توسط یکی از دستورات داخلی پوسته (\en{Shell Builtin}) اشغال شده است. بنابراین، نام دیگری نظیر \cmd{foo} را بررسی می‌کنیم:

\begin{latin}
\begin{lstlisting}
[me@linuxbox ~]$ type foo
bash: type: foo: not found
\end{lstlisting}
\end{latin}

عدم یافتن نام \cmd{foo} نشان‌دهنده‌ی آزاد بودن آن است. اکنون می‌توان نام مستعار را با استفاده از نحوِ \cmd{alias name='string'} تعریف کرد:

\begin{latin}
\begin{lstlisting}
[me@linuxbox ~]$ alias foo='cd /usr; ls; cd -'
\end{lstlisting}
\end{latin}

با اجرای این فرمان، نام مستعار \cmd{foo} به توالیِ سه‌گانه‌ی دستورات مذکور اختصاص می‌یابد. از این پس، در هر بخشی از محیط خط فرمان که نیاز به اجرای آن فرآیند باشد، صرفاً با تایپ \cmd{foo} عملیات اجرا می‌گردد:

\begin{latin}
\begin{lstlisting}
[me@linuxbox ~]$ foo
bin  games  include  lib  local  sbin  share  src
/home/me
[me@linuxbox ~]$
\end{lstlisting}
\end{latin}

برای اعتبارسنجی و اطمینان از صحت انتساب، مجدداً از \cmd{type} استفاده می‌کنیم:

\begin{latin}
\begin{lstlisting}
[me@linuxbox ~]$ type foo
foo is aliased to `cd /usr; ls; cd -'
\end{lstlisting}
\end{latin}

در صورتی که دیگر به این نام مستعار نیاز نداشته باشید، می‌توانید با استفاده از دستور \cmd{unalias} آن را از حافظه‌ی نشست جاری حذف کنید:

\begin{latin}
\begin{lstlisting}
[me@linuxbox ~]$ unalias foo
[me@linuxbox ~]$ type foo
bash: type: foo: not found
\end{lstlisting}
\end{latin}

اگرچه در مثال پیشین از یک نام جدید برای ایجاد \cmd{alias} استفاده کردیم، اما در دنیای واقعی، نام‌های مستعار اغلب برای افزودن گزینه‌های پیش‌فرض به فرامین موجود به‌کار می‌روند تا رفتار آن‌ها بهینه‌سازی شود. برای نمونه، در بسیاری از توزیع‌های لینوکس، فرمان \cmd{ls} به‌صورت پیش‌فرض یک \cmd{alias} است که قابلیت نمایش رنگی خروجی را فعال می‌کند:

\begin{latin}
\begin{lstlisting}
[me@linuxbox ~]$ type ls
ls is aliased to `ls --color=tty'
\end{lstlisting}
\end{latin}

برای مشاهده فهرست کامل تمامی نام‌های مستعار تعریف‌شده در نشست (\en{Session}) فعلی، کافی است دستور \cmd{alias} را بدون هیچ آرگومان اضافه‌ای اجرا کنید:

\begin{latin}
\begin{lstlisting}
[me@linuxbox ~]$ alias
alias l.='ls -d .* --color=tty'
alias ll='ls -l --color=tty'
alias ls='ls --color=tty'
\end{lstlisting}
\end{latin}

این تنظیمات پیش‌فرض، با هدف ارتقای تجربه کاربری و تسهیل تعامل با رابط خط فرمان (\en{CLI}) طراحی شده‌اند.

\begin{note}
\textbf{نکته:} تعریف نام‌های مستعار به‌صورت مستقیم در خط فرمان، یک محدودیت اساسی دارد: این تعاریف موقت هستند و با بستن پنجره ترمینال یا خاتمه یافتن نشستِ پوسته، از حافظه پاک می‌شوند. در مباحث آینده، نحوه ثبت دائمی این نام‌های مستعار در فایل‌های پیکربندی پوسته (نظیر \file{.bashrc}) را خواهیم آموخت تا در هر بار ورود به سیستم، به‌صورت خودکار بارگذاری شوند.
\end{note}

\section{جمع‌بندی}

اتخاذ این رویکردها، موثرترین راهکار برای گذار از یک کاربر عادی به یک متخصص حرفه‌ای لینوکس است. برای تسلط عمیق بر فرامین، هیچ جایگزینی برای مطالعه‌ی مستندات رسمی و آزمون‌وخطای گزینه‌های مختلف در یک محیط ایزوله و امن وجود ندارد.

اکنون که با ابزارهای مرجع نظیر \cmd{man} و \cmd{info} آشنا شده‌اید، پیشنهاد می‌شود تمامی دستوراتی را که پیش‌تر فرا گرفته‌اید (مانند \cmd{ls}، \cmd{cp}، \cmd{mv}، \cmd{rm} و \cmd{mkdir}) مجدداً مورد بررسی قرار دهید:

\begin{enumerate}
  \item با بهره‌گیری از \cmd{man} یا \cmd{info}، صفحات راهنمای اختصاصی هر دستور را با دقت مطالعه کنید.
  \item به دنبال گزینه‌ها (\en{Options}) و پارامترهای جدیدی بگردید که تاکنون از آن‌ها استفاده نکرده‌اید.
  \item در محیط تمرینی خود (\en{Playground})، این گزینه‌ها را اجرا کنید تا خروجی و رفتار عملیاتی آن‌ها را از نزدیک مشاهده نمایید.
\end{enumerate}

با تکیه بر این متدولوژی، درک فنی شما از منطق فرامین به‌طور چشمگیری ارتقا یافته و توانمندی شما در مدیریت کارآمدتر و ایمن‌تر سیستم‌عامل لینوکس تثبیت خواهد شد.

\section{منابع مطالعه تکمیلی}

علاوه بر مستندات داخلی سیستم، منابع آنلاین متعددی برای تعمیق دانش در زمینه لینوکس و رابط خط فرمان (\en{CLI}) در دسترس است. در ادامه، برخی از معتبرترین مراجع تخصصی معرفی می‌شوند:

راهنمای مرجع \en{Bash} یک راهنمای جامع برای پوسته \en{Bash} است. این منبع با وجود اینکه یک اثر مرجع محسوب می‌شود، حاوی مثال‌های کاربردی بوده و مطالعه آن نسبت به صفحات راهنمای داخلی \en{Bash} بسیار آسان‌تر است.

\begin{latin}
\url{http://www.gnu.org/software/bash/manual/bashref.html}
\end{latin}

مستند پرسش‌های متداول \en{Bash} شامل پاسخ به سوالات پرتکرار پیرامون این پوسته است. این فهرست اگرچه برای کاربران سطح متوسط تا پیشرفته تدوین شده، اما حاوی اطلاعات بسیار ارزشمندی است.

\begin{latin}
\url{http://mywiki.wooledge.org/BashFAQ}
\end{latin}

پروژه گنو (\en{GNU Project}) مستندات گسترده‌ای را برای برنامه‌های خود ارائه می‌دهد که هسته اصلی تجربه کار با خط فرمان لینوکس را تشکیل می‌دهند. فهرست کامل این مستندات در اینجا قابل مشاهده است:

\begin{latin}
\url{http://www.gnu.org/manual/manual.html}
\end{latin}

ویکی‌پدیا نیز مقاله جالبی درباره صفحات راهنما (\en{man pages}) دارد:

\begin{latin}
\url{http://en.wikipedia.org/wiki/Man_page}
\end{latin}